\documentclass[10pt,a4paper]{amsart}
\usepackage[margin=19mm]{geometry}
\usepackage{amsmath,amssymb,mathtools}
\usepackage[scr=rsfs,cal=euler]{mathalfa}
\usepackage{xfrac}
\usepackage[unicode,hidelinks,bookmarks=false]{hyperref}
\newcommand{\ie}{i.e.\ }
\newcommand{\cf}{cf.\ }
\newcommand{\resp}{respectively\ }
\newcommand{\Subsection}{\subsection}
\providecommand{\nobreakdash}{\nobreak\hbox{-}}
\setlength{\emergencystretch}{2em}
\multlinegap0pt
\usepackage{amssymb}
\usepackage{mathtools}
\newtheorem{proposition}{Proposition}
\newtheorem{lemma}{Lemma}
\newtheorem{theorem}{Theorem}
\newcommand{\e}{\ensuremath{\mathrm{e}}}
\title{Error bounds for splitting methods in unitary problems}
\author{Fernando Casas, Ander Murua}
\date{Source: arXiv:2604.01026v1 (CC BY 4.0)}
\begin{document}
\maketitle

\noindent\emph{Source and licence.} Error bounds for splitting methods in unitary problems. Version arXiv:2604.01026v1 (CC BY 4.0), distributed by arXiv under the Creative Commons Attribution 4.0 International licence, http://creativecommons.org/licenses/by/4.0/, as recorded in the arXiv licence metadata for this exact version and shown on the abstract page https://arxiv.org/abs/2604.01026v1. Full licence notice: \texttt{licenses/arxiv-2604.01026-CC-BY-4.0.txt}.\par\smallskip

\noindent\textit{Editorial scope.} Counted unit: the article's theorem th\_compo (source0.tex lines 481--495). The article's complete original proof of it is delivered with it (source0.tex lines 497--512, one \texttt{proof} environment of 16 lines). No mathematics is written, repaired or re-derived: the statement and the proof are the article's own lines, in the article's own notation, and no retained line is substituted or re-flowed.  Retained uncounted as the source-local context the proof needs: lines 248--259; lines 427--430; lines 456--466.  Nothing was omitted from any retained span, so the package carries no omission record.  No retained line contains a citation, so the wrapper carries no bibliography and no bibliography entry.  No retained line contains a figure environment, an image-inclusion command or drawing code, so no image is delivered and the package declares no asset dependency.  Editorial formatting only, in addition: an AMS \texttt{amsart} preamble that restates the article's own declarations and macros used by the retained lines, namely the environments \texttt{proposition}, \texttt{lemma}, \texttt{theorem} on separate counters, together with the standard packages those lines need.  The retained environments print in this document as Proposition 1, Lemma 1, Theorem 1 in order of appearance, whereas the article prints the same items with the numbering of its own full document.  The complete native source of the article as distributed in the arXiv e-print for this exact version is bundled unchanged as \texttt{sources/197630/source0.tex}.
\begin{proposition}
\label{p:boundsB}
If $C_1, \ldots, C_r$ are bounded skew-adjoint operators, then, for all $q \geq 1$,
\begin{equation}
\label{eq:expB2}
\mathrm{e}^{h C_r} \cdots \mathrm{e}^{h C_1} = I + \sum_{n=1}^{q} h^n \sum_{j_1+\cdots+j_r = n} \frac{C_r^{j_r} \cdots C_1^{j_1}}{j_r! \cdots j_1!} + h^{q+1} \mathcal{R}_{q+1}(h),
\end{equation}
where
\begin{equation*}
\|\mathcal{R}_{q+1}(h)\| \leq \frac{\big(\|C_1\| + \cdots + \|C_r\| \big)^{q+1}}{(q+1)!}.
\end{equation*}
\end{proposition}

\begin{equation} \label{palindromic_alpha}
  \Psi(h) = \chi(\alpha_{2s} h) \, \chi^*(\alpha_{2s-1} h) \cdots \chi(\alpha_2 h) \, \chi^*(\alpha_1 h), 
     \quad \mbox{ with } \quad \alpha_{2s+1-j} = \alpha_j, \quad 1 \leq j \leq 2s
\end{equation}

\begin{lemma} \label{l:Theta}
Let  $\Phi(h)$ be a one-parameter family of unitary operators, and consider the symmetric factorization $\Psi(h) = \Phi^*(h) \, \Phi(h)$.
Then
\begin{equation*}
\|\Psi(h) - \e^{h\,  (A_1+\cdots+A_N)}\| \leq \| \Theta(h) - \Theta(-h)\|,
\end{equation*}
where
\begin{equation} \label{eq:Theta}
\Theta(h) = \Phi(h) \, \e^{-\frac{h}{2}  (A_1+\cdots+A_N)}.
\end{equation}
\end{lemma}

\begin{theorem} \label{th_compo}
Suppose that the time-symmetric $2s$-stage composition \eqref{palindromic_alpha} 
\[
   \Psi(h) = \chi(\alpha_{1} h) \, \chi^*(\alpha_{2} h) \cdots \chi(\alpha_2 h) \, \chi^*(\alpha_1 h), 
\]
with $(\alpha_1, \alpha_2, \ldots, \alpha_s) \in \mathbb{R}^s$ and
\[
  \chi(h) = \e^{h A_1} \, \e^{h A_2} \cdots \e^{h A_N}, \qquad \chi^*(h) = \e^{h A_N} \cdots \e^{h A_2} \, \e^{h A_1}, 
\]
 provides an approximation of even order $p$ to $\e^{h(A_1 + \cdots + A_N)}$,
with $A_j$ skew-adjoint operators verifying that $\|A_j\| = a_j < \infty$ for $j=1,\ldots, N$. Then, the following estimate holds:
\begin{equation*} 
  \left\|\Psi(h)-\e^{h \sum_{j=1}^N A_j} \right\| \le \frac{2^{-p}}{(p+1)!} \left( 1+ 2 \sum_{i=1}^s |\alpha_i| \right)^{p+1}  \, \left(h \sum_{j=1}^N a_j \right)^{p+1}.
\end{equation*}
\end{theorem}

\begin{proof}
It is straightforward to check that 
\[
\Theta(h)=\e^{h C_r} \cdots \e^{h C_1}, \qquad 
\Theta(-h) = \e^{-h C_r} \cdots \e^{-h C_1}
\]
 for certain bounded skew-adjoint operators $C_j$ such that 
\begin{equation*}
\|C_1\|+\cdots+\|C_r\|  = \left( \frac{1}{2} + |\alpha_1|+\cdots+|\alpha_s| \right) \, \sum_{j=1}^N a_j.
\end{equation*}
Since $\mathcal{E}(h) = \mathcal{O}(h^{p+1})$, Lemma~\ref{l:Theta} and Proposition~\ref{p:boundsB} lead to
\[
  \|\mathcal{E}(h)\| \le  \frac{2 h^{p+1}}{(p+1)!} \left( \|C_1\|+\cdots+\|C_r\| \right)^{p+1}
 \]
 whence the conclusion follows. 
\end{proof}

\end{document}
